% !TEX program = pdflatex
\documentclass[10pt]{article}

\usepackage[utf8]{inputenc}
\usepackage[T1]{fontenc}
\usepackage[english]{babel}
\usepackage{lmodern}
\usepackage{microtype}
\usepackage[margin=0.72in]{geometry}
\usepackage{amsmath,amssymb,amsthm,mathtools}
\usepackage{booktabs}
\usepackage{xcolor}
\usepackage{enumitem}
\usepackage[most]{tcolorbox}
\usepackage{fancyhdr}
\usepackage[hidelinks]{hyperref}

\definecolor{CourseBlue}{HTML}{1A3C68}
\definecolor{CourseGreen}{HTML}{216E39}
\definecolor{CourseGold}{HTML}{B8860B}
\definecolor{SoftBlue}{HTML}{F2F6FA}
\definecolor{SoftGreen}{HTML}{F1F8F3}
\definecolor{SoftGold}{HTML}{FFF8E8}

\setlength{\parindent}{0pt}
\setlength{\parskip}{0.32em}
\setlist[itemize]{leftmargin=1.5em,itemsep=0.12em,topsep=0.2em}
\setlist[enumerate]{leftmargin=1.7em,itemsep=0.12em,topsep=0.2em}
\allowdisplaybreaks

\pagestyle{fancy}
\fancyhf{}
\lhead{Stochastic Calculus and Option Pricing}
\rhead{Lecture 2 Notes}
\cfoot{\thepage}
\renewcommand{\headrulewidth}{0.3pt}

\newcommand{\E}{\mathbb E}
\newcommand{\Pbb}{\mathbb P}
\newcommand{\Var}{\operatorname{Var}}
\newcommand{\Cov}{\operatorname{Cov}}
\newcommand{\R}{\mathbb R}
\newcommand{\1}{\mathbf 1}
\newcommand{\cF}{\mathcal F}
\newcommand{\dd}{\,\mathrm d}

\theoremstyle{definition}
\newtheorem{definition}{Definition}[section]
\newtheorem{proposition}[definition]{Proposition}
\newtheorem{theorem}[definition]{Theorem}
\newtheorem{corollary}[definition]{Corollary}
\newtheorem{remark}[definition]{Remark}

\newtcolorbox{boardexample}[1]{
  enhanced,breakable,colback=SoftBlue,colframe=CourseBlue,
  boxrule=0.55pt,arc=1mm,left=1.5mm,right=1.5mm,top=1mm,bottom=1mm,
  before skip=0.55em,after skip=0.55em,
  title={Example --- #1},fonttitle=\bfseries}
\newtcolorbox{studentexercise}[1]{
  enhanced,breakable,colback=SoftGold,colframe=CourseGold,
  boxrule=0.55pt,arc=1mm,left=1.5mm,right=1.5mm,top=1mm,bottom=1mm,
  before skip=0.55em,after skip=0.35em,
  title={Exercise --- #1},fonttitle=\bfseries}
\newtcolorbox{shortsolution}[1][]{
  enhanced,breakable,colback=SoftGreen,colframe=CourseGreen,
  boxrule=0.55pt,arc=1mm,left=1.5mm,right=1.5mm,top=1mm,bottom=1mm,
  before skip=0.25em,after skip=0.6em,
  title={Short solution#1},fonttitle=\bfseries}
\newtcolorbox{keyidea}[1]{
  enhanced,breakable,colback=white,colframe=CourseGreen,
  boxrule=0.7pt,arc=1mm,left=1.5mm,right=1.5mm,top=1mm,bottom=1mm,
  before skip=0.5em,after skip=0.5em,
  title={#1},fonttitle=\bfseries}

\title{\textbf{Brownian Motion, Stopping Times, and Quadratic Variation}\\
\large Lecture 2 Notes --- Class Support and Student Summary}
\author{}
\date{}

\begin{document}
\maketitle
\vspace{-1.4em}

% ============================================================
\section{From Random Walks to Brownian Motion}

\subsection{Discrete intuition and scaling}

Let \(\varepsilon_1,\varepsilon_2,\ldots\) be independent with
\(\Pbb(\varepsilon_k=1)=\Pbb(\varepsilon_k=-1)=1/2\), and set
\[
S_n=\sum_{k=1}^n\varepsilon_k.
\]
Then \(\E[S_n]=0\) and \(\Var(S_n)=n\). A typical displacement therefore
has order \(\sqrt n\), not order \(n\).

\begin{proposition}[Mean and variance of the random walk]
For every \(n\geq0\),
\[
\E[S_n]=0,
\qquad
\Var(S_n)=n.
\]
\end{proposition}

\begin{proof}
Linearity gives \(\E[S_n]=\sum_k\E[\varepsilon_k]=0\). Independence gives
\(\Var(S_n)=\sum_k\Var(\varepsilon_k)=n\).
\end{proof}

To fit \(n\) shocks into one unit of time while keeping a non-degenerate
variance, define
\[
W_t^{(n)}=\frac{S_{\lfloor nt\rfloor}}{\sqrt n},
\]
with linear interpolation between grid points. For grid times \(s<t\),
\[
\Var(W_t^{(n)}-W_s^{(n)})=t-s.
\]

\begin{theorem}[Donsker's invariance principle, informal]
On every bounded time interval, the rescaled random-walk path \(W^{(n)}\)
converges in distribution to a standard Brownian motion as \(n\to\infty\).
\end{theorem}

This is the functional counterpart of the central limit theorem: many small,
centered, independent shocks produce a Gaussian continuous-time limit.

\begin{studentexercise}{why the square-root normalization is unique}
Suppose each of the \(n\) random-walk shocks is multiplied by \(n^{-\alpha}\).
Find the variance after one unit of time. For which \(\alpha\) does it converge
to a finite nonzero limit?
\end{studentexercise}
\begin{shortsolution}
The variance is \(n(n^{-\alpha})^2=n^{1-2\alpha}\). It has a finite nonzero
limit only when \(1-2\alpha=0\), hence \(\alpha=1/2\).
\end{shortsolution}

% ============================================================
\section{Brownian Motion and Gaussian Structure}

\subsection{Definition and finite-dimensional laws}

\begin{definition}[Standard Brownian motion]
A process \(W=(W_t)_{t\geq0}\) is a standard Brownian motion if:
\begin{enumerate}[label=(BM\arabic*)]
  \item \(W_0=0\) almost surely;
  \item increments over disjoint intervals are independent;
  \item \(W_t-W_s\sim\mathcal N(0,t-s)\) for \(0\leq s<t\);
  \item its paths are continuous almost surely.
\end{enumerate}
\end{definition}

Continuity means no jumps. It does not imply smoothness: Brownian paths are
extremely irregular at every time scale.

\begin{definition}[Gaussian vector and Gaussian process]
A vector \(Z\in\R^m\) is Gaussian if \(a^\top Z\) is Gaussian for every
\(a\in\R^m\). A process \(X\) is Gaussian if
\((X_{t_1},\ldots,X_{t_m})\) is Gaussian for every finite collection of
times.
\end{definition}

\begin{proposition}[Brownian motion is a centered Gaussian process]
For every \(0<t_1<\cdots<t_m\), the vector
\((W_{t_1},\ldots,W_{t_m})\) is centered Gaussian.
\end{proposition}

\begin{proof}
Let \(\Delta_j=W_{t_j}-W_{t_{j-1}}\), with \(t_0=0\). The vector
\((\Delta_1,\ldots,\Delta_m)\) has independent Gaussian coordinates, and
\[
W_{t_k}=\sum_{j=1}^k\Delta_j.
\]
The level vector is therefore a linear transformation of a Gaussian vector.
\end{proof}

\begin{proposition}[Brownian covariance kernel]
For \(s,t\geq0\),
\[
\E[W_t]=0,
\qquad
K(s,t):=\Cov(W_s,W_t)=\min(s,t).
\]
The kernel \(K\) is positive semidefinite.
\end{proposition}

\begin{proof}
If \(s\leq t\), write \(W_t=W_s+(W_t-W_s)\). The second term is centered
and independent of \(W_s\), so
\(\Cov(W_s,W_t)=\Var(W_s)=s\). Moreover,
\[
\sum_{i,j}a_ia_jK(t_i,t_j)
=\Var\!\left(\sum_i a_iW_{t_i}\right)\geq0.
\]
\end{proof}

A Gaussian law is determined by its mean vector and covariance matrix.
Consequently, the zero mean and the kernel \(K(s,t)=\min(s,t)\) determine all
finite-dimensional Brownian laws.

\begin{boardexample}{a symbolic covariance matrix}
For \(0<r<s<t\),
\[
\Cov\!\left((W_r,W_s,W_t)^\top\right)
=\begin{pmatrix}
r&r&r\\ r&s&s\\ r&s&t
\end{pmatrix}.
\]
The repeated entries represent shared history.
\end{boardexample}

\subsection{Transformations and path properties}

\begin{proposition}[Reusable Brownian transformations]
Fix \(c>0\) and \(s>0\). The processes
\[
(-W_t)_{t\geq0},
\qquad
\left(\frac{W_{ct}}{\sqrt c}\right)_{t\geq0},
\qquad
(W_{s+t}-W_s)_{t\geq0}
\]
are standard Brownian motions. Moreover,
\[
(W_{s-t}-W_s)_{0\leq t\leq s}
\]
has the law of Brownian motion restricted to \([0,s]\).
\end{proposition}

\begin{proof}
Each process starts from zero and is continuous. Its increments are linear
transforms of Brownian increments over disjoint intervals, hence remain
independent and Gaussian. Their variances are respectively \(t-u\),
\((c(t-u))/c=t-u\), \(t-u\), and \(t-u\). These are the four defining
properties.
\end{proof}

The scaling property gives the square-root-of-time rule:
\[
W_{ct}\overset d=\sqrt c\,W_t.
\]

\begin{theorem}[Brownian roughness]
Almost surely, a Brownian path is nowhere differentiable and has infinite
total variation on every nontrivial interval.
\end{theorem}

The useful heuristic is
\[
\frac{W_{t+h}-W_t}{h}\sim\mathcal N(0,1/h).
\]
The variance of the difference quotient explodes as \(h\downarrow0\), so a
stable tangent is impossible. These pathwise theorems are stated without
proof here.

\begin{proposition}[Oscillation and the law of the iterated logarithm]
Almost surely,
\[
\limsup_{t\to\infty}W_t=+\infty,
\qquad
\liminf_{t\to\infty}W_t=-\infty,
\]
and
\[
\begin{aligned}
\limsup_{t\to\infty}\frac{W_t}{\sqrt{2t\log\log t}}&=1,
&\liminf_{t\to\infty}\frac{W_t}{\sqrt{2t\log\log t}}&=-1,\\
\limsup_{t\downarrow0}\frac{W_t}{\sqrt{2t\log\log(1/t)}}&=1,
&\liminf_{t\downarrow0}\frac{W_t}{\sqrt{2t\log\log(1/t)}}&=-1.
\end{aligned}
\]
\end{proposition}

\begin{studentexercise}{time reversal}
Fix \(s>0\) and set \(R_t=W_{s-t}-W_s\) for \(0\leq t\leq s\). Verify
directly the distribution of \(R_t-R_u\) for \(u<t\), and explain why
increments over disjoint intervals are independent.
\end{studentexercise}
\begin{shortsolution}
\[
R_t-R_u=W_{s-t}-W_{s-u}\sim\mathcal N(0,t-u).
\]
Disjoint intervals in reversed time correspond to disjoint intervals of the
original Brownian path. Continuity is inherited, so \(R\) has Brownian law on
\([0,s]\).
\end{shortsolution}

% ============================================================
\section{Markov Renewal and Martingales}

Let \(\cF_t=\sigma(W_u:0\leq u\leq t)\), completed in the usual way. It
represents the entire observed Brownian path up to time \(t\).

\begin{proposition}[Deterministic-time renewal]
For fixed \(t\geq0\),
\[
\widetilde W_u=W_{t+u}-W_t,
\qquad u\geq0,
\]
is a standard Brownian motion independent of \(\cF_t\).
\end{proposition}

\begin{proof}
The Brownian properties follow from stationary independent increments and
continuity. Every increment of \(\widetilde W\) occurs after time \(t\), so
it is independent of the past \(\cF_t\).
\end{proof}

\begin{proposition}[Markov property in conditional-expectation form]
For every bounded measurable \(f\) and \(h>0\),
\[
\boxed{
\E[f(W_{t+h})\mid\cF_t]
=\E[f(W_{t+h})\mid W_t]
=(P_hf)(W_t),
}
\]
where
\[
(P_hf)(x)=\int_{\R}f(x+y)
\frac{e^{-y^2/(2h)}}{\sqrt{2\pi h}}\dd y.
\]
\end{proposition}

\begin{proof}
Write \(W_{t+h}=W_t+\widetilde W_h\), where
\(\widetilde W_h\sim\mathcal N(0,h)\) is independent of \(\cF_t\). When
conditioning on \(\cF_t\), \(W_t\) is known and the remaining expectation is
the displayed Gaussian integral.
\end{proof}

The operators \((P_h)_{h\geq0}\) form a semigroup:
\[
P_{h+k}=P_hP_k.
\]
This follows either from Gaussian convolution or by applying the Markov
property twice.

\begin{proposition}[Brownian motion is a martingale]
For \(0\leq s<t\),
\[
\E[W_t\mid\cF_s]=W_s.
\]
\end{proposition}

\begin{proof}
Use \(W_t=W_s+(W_t-W_s)\). The increment is centered and independent of
\(\cF_s\).
\end{proof}

Markov and martingale are different statements. Markov says that the current
state summarizes the past for the entire future law. Martingale says only that
the current value is the conditional mean of a future value.

\begin{boardexample}{a conditional Gaussian moment}
For \(\theta\in\R\), renewal gives
\[
\E[e^{\theta W_{t+h}}\mid\cF_t]
=e^{\theta W_t}\E[e^{\theta\widetilde W_h}]
=\exp\!\left(\theta W_t+\frac12\theta^2h\right).
\]
\end{boardexample}

\begin{studentexercise}{conditional second moment}
For \(0\leq s<t\), compute \(\E[W_t^2\mid\cF_s]\). Deduce a second
Brownian martingale.
\end{studentexercise}
\begin{shortsolution}
Expanding \(W_t=W_s+(W_t-W_s)\) gives
\[
\E[W_t^2\mid\cF_s]=W_s^2+(t-s).
\]
Hence \(\E[W_t^2-t\mid\cF_s]=W_s^2-s\), so \(W_t^2-t\) is a martingale.
\end{shortsolution}

% ============================================================
\section{Stopping Times and Strong Markov Property}

\subsection{No look-ahead}

In discrete time, a random integer \(\tau\) is a stopping time when
\(\{\tau\leq n\}\in\cF_n\) for every \(n\). A first hit can be detected as
the path unfolds; a last hit before a future horizon cannot.

\begin{definition}[Stopping time]
A random time \(\tau:\Omega\to[0,\infty]\) is a stopping time for
\((\cF_t)\) if
\[
\{\tau\leq t\}\in\cF_t
\qquad\text{for every }t\geq0.
\]
\end{definition}

The event \(\{\tau\leq t\}\) asks whether the rule has already triggered.
The answer may use the observed path through time \(t\), but not the future.

\begin{proposition}[Canonical examples]
The following are stopping times:
\[
t_0,
\qquad
\tau_a=\inf\{t\geq0:W_t=a\},
\qquad
\tau_{\ell,u}=\inf\{t\geq0:W_t\notin(\ell,u)\}.
\]
For fixed \(T>0\), the last zero
\(\rho_T=\sup\{t\leq T:W_t=0\}\) and the time of the maximum on \([0,T]\)
are not stopping times.
\end{proposition}

\begin{proof}
The deterministic time is immediate. By continuity,
\[
\{\tau_a\leq t\}
=\left\{\max_{0\leq v\leq t}W_v\geq a\right\}
\]
for \(a>0\), and the right-hand event is observable by time \(t\). The same
argument works for first exit. Identifying a last zero or the global maximum
requires checking that no later event occurs, which uses future information.
\end{proof}

\begin{definition}[Information available at a stopping time]
For a stopping time \(\tau\),
\[
\cF_\tau
=\{A\in\cF:A\cap\{\tau\leq t\}\in\cF_t
\text{ for every }t\geq0\}.
\]
\end{definition}

An event in \(\cF_\tau\) can be decided using only information revealed by
the moment the rule triggers. For continuous adapted \(W\) and almost surely
finite \(\tau\), the stopped value \(W_\tau\) is \(\cF_\tau\)-measurable.

\begin{theorem}[Strong Markov property]
If \(\tau\) is an almost surely finite stopping time, then
\[
(W_{\tau+u}-W_\tau)_{u\geq0}
\]
is a standard Brownian motion independent of \(\cF_\tau\).
\end{theorem}

\begin{remark}[Proof idea]
Round \(\tau\) upward to dyadic stopping times
\(\tau_n=2^{-n}\lceil2^n\tau\rceil\). On each event
\(\{\tau_n=k2^{-n}\}\), apply deterministic renewal. Then let
\(n\to\infty\) and use path continuity. The full measure-theoretic proof is
not needed in this lecture.
\end{remark}

\begin{studentexercise}{first hit after a deterministic date}
Fix \(s>0\) and define
\[
\tau=\inf\{t\geq s:W_t=0\}.
\]
Is \(\tau\) a stopping time? Contrast it with the last zero before a fixed
horizon \(T\).
\end{studentexercise}
\begin{shortsolution}
Yes. By time \(t\), one can inspect the observed path on \([s,t]\) and decide
whether a zero occurred. The last zero before \(T\) is not a stopping time,
because at time \(t<T\) one cannot rule out a later return to zero.
\end{shortsolution}

% ============================================================
\section{Reflection Principle and Hitting Probabilities}

Let
\[
M_t=\max_{0\leq u\leq t}W_u,
\qquad
\tau_a=\inf\{u\geq0:W_u=a\},
\qquad a>0.
\]
By continuity, \(\{M_t\geq a\}=\{\tau_a\leq t\}\).

On the event \(\{\tau_a\leq t\}\), define the reflected path
\[
W_u^*=\begin{cases}
W_u,&u\leq\tau_a,\\
2a-W_u,&u>\tau_a.
\end{cases}
\]
Strong Markov says that the continuation after \(\tau_a\) is fresh Brownian
motion; symmetry says that its negative has the same law. Therefore reflection
preserves Brownian probability.

\begin{theorem}[Reflection principle]
For \(a>0\) and \(t>0\),
\[
\boxed{
\Pbb(M_t\geq a)=2\Pbb(W_t\geq a)
=2\left(1-\Phi\!\left(\frac{a}{\sqrt t}\right)\right).
}
\]
\end{theorem}

\begin{proof}
Split \(\{M_t\geq a\}\) according to the terminal value:
\[
\{M_t\geq a\}
=\{W_t\geq a\}\,\dot\cup\,
\{M_t\geq a,W_t<a\}.
\]
Reflection after \(\tau_a\) is a bijection from the second event to
\(\{W_t>a\}\) and preserves probability. Since \(\Pbb(W_t=a)=0\), the two
pieces have the same probability.
\end{proof}

\begin{corollary}[First-passage distribution and density]
For \(a>0\),
\[
\Pbb(\tau_a\leq t)
=2\left(1-\Phi\!\left(\frac{a}{\sqrt t}\right)\right),
\]
and
\[
f_{\tau_a}(t)
=\frac{a}{\sqrt{2\pi t^3}}
\exp\!\left(-\frac{a^2}{2t}\right),
\qquad t>0.
\]
\end{corollary}

\begin{corollary}[Maximum and scaling of the hitting time]
\[
M_t\overset d=|W_t|,
\qquad
\E[M_t]=\sqrt{\frac{2t}{\pi}},
\qquad
\tau_a\overset d=a^2\tau_1.
\]
Moreover, \(\Pbb(\tau_a<\infty)=1\), while \(\E[\tau_a]=\infty\).
\end{corollary}

The infinite mean follows from
\[
\Pbb(\tau_a>t)
=2\Phi(a/\sqrt t)-1
\sim a\sqrt{\frac{2}{\pi t}},
\]
whose integral diverges at infinity.

\begin{boardexample}{a symbolic barrier probability}
Suppose a driftless log-price satisfies
\(\log(S_t/S_0)=\sigma W_t\), and let \(H>S_0\). Then
\[
\Pbb\!\left(\max_{0\leq u\leq T}S_u\geq H\right)
=2\left[1-\Phi\!\left(
\frac{\log(H/S_0)}{\sigma\sqrt T}
\right)\right].
\]
Reflection converts a path-dependent barrier event into a Gaussian tail.
\end{boardexample}

\begin{studentexercise}{distribution and mean of the maximum}
Derive \(\Pbb(M_t\leq a)\) for \(a\geq0\), and compute \(\E[M_t]\).
\end{studentexercise}
\begin{shortsolution}
\[
\Pbb(M_t\leq a)=2\Phi(a/\sqrt t)-1.
\]
Thus \(M_t\overset d=|W_t|\), and the mean of a half-normal variable is
\(\E[M_t]=\sqrt{2t/\pi}\).
\end{shortsolution}

% ============================================================
\section{Quadratic Variation and Brownian Roughness}

\subsection{Definition and basic properties}

For a partition \(\pi=\{0=t_0<\cdots<t_n=t\}\), let
\(|\pi|=\max_i(t_{i+1}-t_i)\).

\begin{definition}[Quadratic variation]
For a continuous process \(X\), define
\[
[X]_t
=\lim_{|\pi|\to0}\sum_i(X_{t_{i+1}}-X_{t_i})^2,
\]
when the limit exists in probability.
\end{definition}

Quadratic variation measures accumulated fine-scale fluctuation. It ignores
the direction of increments and is different from ordinary total variation.

\begin{proposition}[Finite-variation paths have zero quadratic variation]
If \(A\) is continuously differentiable on \([0,T]\), then \([A]_T=0\).
More generally, the result holds for continuous finite-variation paths.
\end{proposition}

\begin{proof}
For \(A\in C^1\), the mean value theorem gives
\[
\sum_i(\Delta_iA)^2
\leq\|A'\|_\infty^2\sum_i(\Delta_it)^2
\leq\|A'\|_\infty^2T|\pi|\longrightarrow0.
\]
\end{proof}

\begin{theorem}[Brownian quadratic variation]
For every \(T\geq0\),
\[
\boxed{[W]_T=T.}
\]
More precisely, for every partition \(\pi\) of \([0,T]\),
\[
\E\!\left[
\left(\sum_i(\Delta_iW)^2-T\right)^2
\right]
\leq2T|\pi|.
\]
\end{theorem}

\begin{proof}
The independent increments satisfy
\(\Delta_iW\sim\mathcal N(0,\Delta_it)\). Hence
\[
\E\!\left[\sum_i(\Delta_iW)^2\right]
=\sum_i\Delta_it=T.
\]
For \(G\sim\mathcal N(0,v)\), \(\Var(G^2)=2v^2\). Therefore
\[
\Var\!\left(\sum_i(\Delta_iW)^2\right)
=2\sum_i(\Delta_it)^2
\leq2|\pi|\sum_i\Delta_it
=2T|\pi|.
\]
The realized sum converges to \(T\) in \(L^2\), hence in probability.
\end{proof}

\begin{proposition}[Drift disappears from quadratic variation]
If \(A\) is continuously differentiable and \(X=A+\sigma W\), then
\[
[X]_t=\sigma^2t.
\]
In particular, \([aX]_t=a^2[X]_t\) and \([X+A]_t=[X]_t\).
\end{proposition}

\begin{proof}
Expand the squared increments. The sum of \((\Delta A)^2\) vanishes. For the
cross-term, Cauchy--Schwarz gives
\[
\left|\sum_i\Delta_iA\,\Delta_iW\right|
\leq
\left(\sum_i(\Delta_iA)^2\right)^{1/2}
\left(\sum_i(\Delta_iW)^2\right)^{1/2},
\]
which tends to zero in probability. The Brownian term tends to
\(\sigma^2t\).
\end{proof}

\begin{center}
\begin{tabular}{lll}
\toprule
Process & Quadratic variation & Fine-scale component \\
\midrule
smooth or finite-variation \(A\) & \([A]_t=0\) & none \\
Brownian motion \(W\) & \([W]_t=t\) & one variance unit per unit time \\
\(A_t+\sigma W_t\) & \([X]_t=\sigma^2t\) & volatility, not drift \\
\bottomrule
\end{tabular}
\end{center}

\subsection{The bracket and the new calculus}

\begin{proposition}[Quadratic-variation martingale]
For a continuous square-integrable martingale \(X\), there is a unique
increasing predictable process \(\langle X\rangle\), starting from zero, such
that
\[
\boxed{X_t^2-\langle X\rangle_t\text{ is a martingale}.}
\]
For continuous martingales, \(\langle X\rangle=[X]\). In particular,
\(\langle W\rangle_t=t\).
\end{proposition}

For Brownian motion this can be checked without general theory:
\[
\E[W_t^2-t\mid\cF_s]
=W_s^2+(t-s)-t
=W_s^2-s.
\]

The It\^o multiplication rules summarize the scale of the increments:
\[
(\dd W_t)^2=\dd t,
\qquad
\dd W_t\,\dd t=0,
\qquad
(\dd t)^2=0.
\]
These are shorthand for limiting statements about quadratic variation, not
ordinary algebraic identities.

\begin{studentexercise}{realized variance with drift}
Let \(X_t=A_t+\sigma W_t\), where \(A\in C^1\), and observe \(X\) on a
partition of \([0,T]\). Show that
\[
\widehat\sigma_\pi^2
=\frac1T\sum_i(X_{t_{i+1}}-X_{t_i})^2
\]
converges in probability to \(\sigma^2\) as \(|\pi|\to0\).
\end{studentexercise}
\begin{shortsolution}
The numerator is the realized quadratic variation of \(X\), which converges
to \([X]_T=\sigma^2T\). Division by \(T\) gives the result.
\end{shortsolution}

\begin{keyidea}{Numerical interpretation}
On observed data, \(\sum_i(\Delta_iX)^2\) is the realized variance. Refining
the grid removes smooth drift and reveals accumulated volatility. In actual
market data, extremely fine grids may be contaminated by bid--ask bounce and
other microstructure noise, so ``finer'' is not automatically ``better.''
\end{keyidea}

% ============================================================
\section{Geometric Brownian Motion}

An additive model \(S_t=S_0+\mu t+\sigma W_t\) can become negative and has
constant absolute volatility. A multiplicative model instead makes log-returns
Gaussian.

\begin{definition}[Geometric Brownian motion]
For \(S_0>0\) and \(\sigma>0\),
\[
S_t=S_0\exp\!\left(
\left(\mu-\frac12\sigma^2\right)t+\sigma W_t
\right).
\]
Then \(S_t>0\) for every \(t\).
\end{definition}

\begin{proposition}[Laws and moments of GBM]
For \(0\leq s<t\),
\[
\log(S_t/S_s)\mid\cF_s
\sim\mathcal N\!\left(
(\mu-\tfrac12\sigma^2)(t-s),\sigma^2(t-s)
\right),
\]
independently of \(\cF_s\). Moreover,
\[
\E[S_t]=S_0e^{\mu t},
\qquad
\Var(S_t)=S_0^2e^{2\mu t}(e^{\sigma^2t}-1),
\]
and
\[
\operatorname{Median}(S_t)
=S_0e^{(\mu-\sigma^2/2)t}.
\]
\end{proposition}

\begin{proof}
The log-return is a deterministic drift plus
\(\sigma(W_t-W_s)\), giving the conditional normal law. For
\(G\sim\mathcal N(0,t)\),
\(\E[e^{\lambda G}]=e^{\lambda^2t/2}\). Applying this with
\(\lambda=\sigma\) and \(2\sigma\) gives the first two moments and hence the
variance. The median is obtained by setting the centered Gaussian term to its
median zero.
\end{proof}

For \(K>0\),
\[
\Pbb(S_t>K)
=1-\Phi\!\left(
\frac{\log(K/S_0)-(\mu-\sigma^2/2)t}{\sigma\sqrt t}
\right).
\]

\begin{proposition}[Quadratic variation of the log-price]
\[
[\log S]_t=\sigma^2t.
\]
\end{proposition}

\begin{proof}
The log-price is a continuously differentiable drift plus \(\sigma W_t\).
The drift has zero quadratic variation and zero limiting cross-term with the
Brownian part.
\end{proof}

GBM is tractable and preserves positivity, but it omits jumps, heavy tails,
changing volatility, and volatility smiles. Brownian motion is a foundational
building block rather than a complete market description.

\begin{keyidea}{What to retain}
Brownian motion is a centered Gaussian process with covariance
\(\min(s,t)\), fresh increments, and continuous rough paths. Markov renewal
describes what happens after a current or stopping time; reflection converts
barrier events into Gaussian tails; quadratic variation records accumulated
volatility and creates the correction terms of stochastic calculus.
\end{keyidea}

\end{document}
